\documentclass[conference,a4paper]{IEEEtran}
\usepackage[T1]{fontenc}
\usepackage[utf8]{inputenc}
\usepackage[brazil]{babel}
\usepackage{lmodern}
\usepackage{textcomp}
\usepackage{array}
\usepackage{pgfplots}
\pgfplotsset{compat=1.18}
\usetikzlibrary{babel}
\setlength{\tabcolsep}{4pt}
\usepackage{booktabs}
\usepackage{xcolor}
\usepackage{url}
\usepackage[hidelinks]{hyperref}
% Trecho que o grupo ainda precisa validar (não remover sem o grupo pedir)
\newcommand{\revisar}[1]{\textcolor{red}{[REVISAR: #1]}}
\newcommand{\us}{\,\textmu s}
\renewcommand{\IEEEkeywordsname}{Palavras-chave}
\newcolumntype{P}[1]{>{\raggedright\arraybackslash}p{#1}}

\begin{document}

\title{TCAM ou algoritmo? Uma revisão crítica do lookup por \textit{longest prefix match} em velocidade, memória, atualização e energia}

\author{\IEEEauthorblockN{Gabriel Trigueiro Melo (15511832), Ana Vitória Chaves Viana (12733869),\\
Diogo dos Santos da Rocha (13673281)}
\IEEEauthorblockA{Escola de Artes, Ciências e Humanidades (EACH), Universidade de São Paulo\\
ACH2026 Redes de Computadores}}

\maketitle

\begin{abstract}
Para repassar um pacote IP, o roteador precisa achar na tabela de repasse o prefixo mais longo que casa com o
endereço de destino (LPM). Há duas famílias de solução: memórias ternárias endereçáveis por conteúdo (TCAM) e
estruturas algorítmicas em SRAM, DRAM ou cache, como Tree Bitmap, DXR, Poptrie e SAIL. Esta revisão analisa sete
artigos publicados entre 2001 e 2018 e compara o que cada um mede em quatro eixos: velocidade de lookup, memória,
custo de atualização e energia. Os números publicados não são comparáveis diretamente. Cada artigo mede com
rigor o eixo em que a sua proposta é forte e trata os demais por modelo ou de forma qualitativa, e grandezas com o mesmo
nome, como memória e custo de atualização, medem coisas diferentes em cada um. Normalizados por cenário, porém,
os resultados deixam de se anular. A disputa entre Poptrie e SAIL, por exemplo, se explica pela fração do tráfego que
casa prefixos maiores que /24, pelo tamanho da cache e por quem implementou o concorrente. Propomos uma tabela de
cenários com a abordagem mais bem sustentada em cada um e o tipo de evidência, e listamos as medições que nenhum
dos artigos fez, em especial o pior caso de atualização durante rajadas e a energia das soluções algorítmicas.
\end{abstract}

\begin{IEEEkeywords}
longest prefix match, TCAM, lookup de endereço IP, tabela de repasse, Poptrie, SAIL, DXR, Tree Bitmap
\end{IEEEkeywords}

% =====================================================================
\section{Introdução}

O plano de dados de um roteador decide a porta de saída de cada pacote consultando a tabela de repasse (FIB).
Desde o CIDR, essa consulta é a busca pelo prefixo mais longo que casa com o endereço de destino
(\textit{longest prefix match}, LPM), e a busca acontece em duas dimensões, o valor e o comprimento do
prefixo~\cite[p.~10]{ruiz01}. A exigência de velocidade é alta: com pacotes de tamanho mínimo, um enlace de
100~GbE pede 148,8 milhões de lookups por segundo (Mlps)~\cite[\S1]{asai15}. A tabela também cresce. Tabelas BGP
completas passavam de 500 mil rotas em 2014~\cite[\S1]{asai15}, e o SAIL cita mais de 700 mil entradas segundo o
RIPE NCC~\cite[\S I-A]{yang18}.

Há duas famílias de solução. A primeira usa TCAM, uma memória que compara a chave com todas as entradas em paralelo
e devolve a primeira que casa~\cite[p.~38]{shah01}. A segunda guarda a tabela numa estrutura de dados (tries
comprimidas, índices diretos, busca binária em faixas) em SRAM, DRAM ou na cache da CPU e paga cada lookup em acessos
à memória~\cite[pp.~13--14]{ruiz01}. Os artigos da segunda família costumam começar dizendo que a TCAM gasta energia
demais e custa caro~\cite[\S1]{zec12}, \cite[\S1]{asai15}, \cite[\S II]{yang18}. Os da primeira a descrevem como
rápida, simples de gerenciar e de custo competitivo~\cite{zane03}.

Cada artigo mede, porém, numa plataforma própria, com tabelas, tráfego e unidades próprios. O Poptrie se declara de
1,37 a 2,62 vezes mais rápido que o SAIL~\cite[\S4.5]{asai15}. Três anos depois, o SAIL se declara de 1,13 a 1,72
vezes mais rápido que o Poptrie~\cite[\S VI-C]{yang18}. Lidos lado a lado, os números parecem se anular.

Esta revisão pergunta se os resultados publicados permitem escolher uma abordagem para LPM, e em quais condições.
Os números não são comparáveis diretamente, porque cada artigo mede com rigor o eixo em que a proposta é forte e trata
os outros por modelo ou de forma qualitativa. Quando normalizados por cenário, porém, deixam de se anular e mostram
onde cada solução vence.

As contribuições são: tabelas que põem lado a lado, com a condição de medida, os números de velocidade, memória e
atualização dos sete artigos; uma análise da divergência entre Poptrie e SAIL; a comparação entre a escala testada em
cada artigo e os limites de formato de cada estrutura; uma tabela de cenários com a abordagem mais bem sustentada em
cada um e o tipo de evidência; e a lista do que uma comparação futura precisaria informar. A Seção~II descreve o método,
a III resume os fundamentos, a IV analisa os quatro eixos e a escala, a V discute os cenários, a VI trata das
limitações e a VII conclui.

% =====================================================================
\section{Método da revisão}

\subsection{Busca e seleção}
A busca exploratória foi feita no IEEE Xplore, na ACM Digital Library e no Google Scholar, com combinações de
``IP lookup'', ``longest prefix match'', ``TCAM update'' e ``TCAM power''.
\revisar{confirmar as bases e as strings que o grupo de fato usou; o log de 31/10 registra só a busca feita pela IA}
Os critérios de inclusão foram:
publicação revisada por pares; tratar de LPM para repasse IP; trazer dados sobre pelo menos um dos quatro eixos; e
haver versão publicada acessível pela rede da USP. A busca, feita com ajuda da IA (Seção~II-C), levantou oito candidatos
e dois de reserva. Um dos oito, sobre atualização de TCAM (He et al., IEEE/ACM ToN 2018), foi descartado depois da
leitura, porque trata de classificação de pacotes multicampo e não de LPM.
\revisar{o PDF do He et al. não está em ep2\_export; anexar o PDF ou retirar o nome} Os dois de reserva não foram lidos.
De três artigos, a primeira versão obtida não era a que citamos. Do DXR e do survey, tínhamos um preprint e um relatório
técnico, com números e numeração diferentes dos da versão publicada, e trocamos os dois pela versão de periódico. Do SAIL,
usamos a versão estendida da IEEE/ACM ToN de 2018 em vez da do SIGCOMM 2014, porque ela é posterior ao Poptrie e
responde às medições dele~\cite[\S VI-C]{yang18}. Com três integrantes, o mínimo exigido é
$\lceil 1{,}5 \times 3 \rceil = 5$ artigos; o corpus final tem sete (Tabela~\ref{tab:corpus}).

\begin{table*}[t]
\centering
\caption{Corpus da revisão e condições de avaliação de cada artigo}
\label{tab:corpus}
\footnotesize
\begin{tabular}{P{2.3cm} P{2.0cm} P{3.0cm} P{3.4cm} P{2.8cm} P{2.8cm}}
\toprule
Artigo & Família & Plataforma & Tabelas & Tráfego & Concorrentes implementados por \\
\midrule
Ruiz-Sánchez et al., 2001~\cite{ruiz01} & survey & Pentium Pro 200 MHz (p.~21) & MaeEast de 1999, 47.113 prefixos (p.~22) & uniforme por prefixo (p.~22) & FreeBSD, autores originais e os próprios (p.~21) \\
Shah e Gupta, 2001~\cite{shah01} & TCAM, atualização & simulação (p.~43) & MAE-EAST e MAE-WEST, 35--43 mil (Tab.~1) & só traços de update de 3 h (p.~43) & os autores (p.~39) \\
Zane et al., 2003~\cite{zane03} & TCAM, energia & simulação em PC de 800 MHz (\S III-D) & 2 reais, 110--122 mil; 1 sintética de $\approx$1M (Tab.~II, \S III-D) & não usa (métrica não depende de tráfego) & só variantes próprias \\
Eatherton et al., 2004~\cite{eath04} & trie multibit & análise e projeto de hardware (\S7) & Mae-East e Mae-West, 19--41 mil; 1M aleatória (Tab.~6) & nenhum & números da literatura (\S3) \\
Zec et al., 2012~\cite{zec12} & índice direto e faixas & AMD FX-8150 com 8 núcleos; Xeon W3530 (\S4.1) & 3 tabelas BGP, 412--421 mil (Tab.~1) & sintético uniforme (\S4.1) & os autores (DIR-24-8) e radix do FreeBSD (\S4.3) \\
Asai e Ohara, 2015~\cite{asai15} & trie comprimida & i7-4770K, 1 núcleo, L3 de 8 MiB (\S4) & 35 reais de 509--532 mil; 4 sintéticas de 756--886 mil (Tab.~1) & aleatório, sequencial, repetido e trace de borda (\S4.2) & os autores, todos (\S4) \\
Yang et al., 2018~\cite{yang18} & bitmaps em níveis & FPGA simulado; i7-3520M com 1 núcleo, L3 de 4 MB; GPU; many-core (\S V) & RIPE, 114--495 mil (Tab.~IV); 12 FIBs de um tier-1 & trace do tier-1, aleatório e por prefixo (\S VI-A) & os autores (4 de 5); código aberto do Poptrie (\S VI-A, \S VI-C) \\
\bottomrule
\end{tabular}
\end{table*}

\subsection{Procedimento de análise}
Cada artigo foi lido inteiro e fichado pelo mesmo roteiro: premissas sobre a FIB, o tráfego, o hardware e a taxa de
atualização; plataforma e cargas de teste; quem implementou os concorrentes; números com unidade, condição e
localizador; limitações admitidas pelos autores e limitações que deduzimos. Separamos o que o artigo afirma do que é
dedução nossa, e no texto as deduções aparecem como tal (``pela nossa leitura'', ``conta nossa''). Quando o texto
extraído do PDF era ambíguo, conferimos o número na página renderizada. A extração automática troca ``\textmu s'' por
``ms'' no DXR~\cite[\S4.4]{zec12} e na tabela de tempos do survey~\cite[Tab.~3]{ruiz01}, um erro de três ordens de
grandeza que só a página renderizada mostra. Gráficos sem rótulo numérico não foram usados como fonte de números. Os localizadores seguem a
paginação impressa da versão publicada de cada artigo.

\subsection{Uso de IA}
O enunciado exige o uso de IA generativa. Usamos o Claude Opus 5.5 (Anthropic) em duas etapas. Na primeira, de 31/10 a
01/11/2026, pelo claude.ai, a IA sugeriu candidatos com link; o grupo baixou os PDFs, e a IA os comparou com o que tinha
afirmado antes. Na segunda, a partir de 16/11/2026, usamos o Claude Code com acesso ao repositório do trabalho e criamos três
agentes: um leitor que faz o fichamento de um artigo (rodamos sete em paralelo), um verificador que confere cada número do
texto contra o PDF e um revisor que lê o artigo pelos critérios do enunciado. A IA escreveu o primeiro rascunho deste texto
a partir dos fichamentos e da estrutura aprovada pelo grupo, e a skill \textit{humanizer} foi usada para tirar vícios de
estilo. As referências vieram dos metadados oficiais dos DOIs, nunca da memória do modelo.

A verificação pegou erros da própria IA: tratar como evidência sobre LPM um artigo de classificação de pacotes;
atribuir ao SAIL uma estatística de distribuição de prefixos que ele cita de outros autores; usar números de um preprint do
DXR; ler ``ms'' onde o PDF diz ``\textmu s''; e citar um limiar de ``880 mil rotas'' que não aparece em nenhum dos dois artigos
envolvidos. \revisar{descrever o que o grupo mudou ao revisar o rascunho e quem revisou cada seção}. O log completo e a
declaração de uso estão nas Partes II e III da entrega.

% =====================================================================
\section{Fundamentos}

\subsection{O problema e a taxonomia}
Com $N$ prefixos de até $W$ bits ($W=32$ no IPv4 e 128 no IPv6), o LPM procura, entre os prefixos que casam com o
endereço, o mais longo~\cite[p.~10]{ruiz01}. O survey de Ruiz-Sánchez et al. organiza as soluções algorítmicas por dois
eixos: a dimensão buscada (o valor ou o comprimento do prefixo) e o modo de busca (linear ou binário). A eles somam-se
técnicas auxiliares como expansão de prefixos, compressão e uso da hierarquia de memória~\cite[pp.~13--14]{ruiz01}. O mesmo
survey fixa os critérios que usamos: velocidade de lookup, custo de atualização, memória e escalabilidade com o número e o
comprimento dos prefixos~\cite[p.~20]{ruiz01}. O survey não trata de TCAM: o termo não aparece no texto, e o hardware
dedicado se resume a uma arquitetura de dois bancos de memória~\cite[p.~16]{ruiz01}. Os dois artigos de TCAM do corpus
cobrem essa lacuna.

O survey compara os esquemas por complexidade assintótica~\cite[Tab.~2]{ruiz01} e mede só o tempo de lookup, num
Pentium Pro de 200~MHz, com uma única tabela de 47.113 prefixos e tráfego uniforme por prefixo~\cite[Tab.~3,
p.~22]{ruiz01}. Atualização e memória ficam no plano assintótico. Da comparação, os autores tiram um compromisso entre
tempo de lookup e tempo de atualização~\cite[p.~22]{ruiz01}. As estruturas mais compactas, como a Lulea, não admitem
atualização incremental, e ``rebuilding the whole structure is the only solution''~\cite[p.~18]{ruiz01}. Otimizar a
estrutura para uma tabela só compensa se ela mudar pouco, ``but this is rarely the case for backbone
routers''~\cite[p.~17]{ruiz01}. Os artigos posteriores herdam esse compromisso: DXR e Poptrie mantêm uma árvore radix
ao lado da estrutura de lookup para reconstruir só a parte afetada por cada atualização~\cite[\S3.4]{zec12},
\cite[\S3.5]{asai15}.

\subsection{TCAM}
Uma TCAM guarda cada prefixo como uma entrada ternária, com bits 0, 1 e ``não importa''. O endereço é comparado com todas
as entradas em paralelo, e um codificador de prioridade escolhe a entrada que casa no menor endereço
físico~\cite[p.~38, Fig.~1]{shah01}. Para que essa primeira entrada seja a do prefixo mais longo, prefixos sobrepostos mais
longos precisam ficar em endereços menores. Por isso inserir um prefixo pode exigir deslocar entradas, até $O(N)$ na
solução ingênua~\cite[p.~37]{shah01}. O outro custo é o consumo. Uma TCAM de alta densidade gasta de 12 a 15~W com todas as
entradas habilitadas para busca~\cite[\S I]{zane03}, e o CoolCAMs modela a potência como proporcional ao número de entradas
pesquisadas~\cite[\S II]{zane03}. Em 2003, um dispositivo de 18~Mbit guardava até 512 mil prefixos IPv4~\cite[\S II]{zane03}.

Shah e Gupta propõem duas formas de reduzir o custo de manter a ordem, ambas com o espaço livre no centro da
TCAM~\cite[pp.~40--41]{shah01}. O PLO\_OPT mantém a ordem por comprimento e move no máximo $L/2$ entradas por operação,
com $L=32$ no IPv4. O CAO\_OPT parte da observação de que só prefixos sobrepostos, isto é, da mesma cadeia entre a raiz
e uma folha da trie, precisam respeitar a ordem; assim, cada prefixo fica a no máximo $D/2$ posições de uma entrada
livre, em que $D$ é o comprimento da maior cadeia~\cite[pp.~40--43]{shah01}. Os dois dependem de um software de gerência
que mantém uma trie auxiliar fora do chip e envia as instruções à TCAM~\cite[p.~38]{shah01}. O CoolCAMs trata da energia:
divide a tabela em grupos (\textit{buckets}) e pesquisa um só por lookup. Na seleção de bits, $k$ bits fixos entre os 16
primeiros do endereço escolhem o grupo, e os prefixos com menos de 16 ou mais de 24 bits, menos de 2\% nas tabelas dos
autores, ficam em blocos pesquisados em todo lookup~\cite[\S III-A]{zane03}. Na arquitetura por trie, uma TCAM de
índice pequena, sempre pesquisada inteira, devolve o grupo, e a partição é feita por dois algoritmos sobre a trie,
\textit{subtree-split} e \textit{postorder-split}~\cite[\S IV, Fig.~5]{zane03}.

\subsection{Estruturas algorítmicas}
As quatro propostas algorítmicas do corpus combinam um índice direto pelos primeiros bits com uma estrutura compacta para
o restante. O Tree Bitmap é uma trie multibit em que cada nó tem dois bitmaps, um para os prefixos guardados no nó e outro
para os filhos, que ficam contíguos. Ele não usa \textit{leaf pushing}, o que limita o custo de uma atualização ao tamanho
de um nó~\cite[\S4]{eath04}. O DXR projeta a tabela em faixas de endereço: os $k$ primeiros bits ($k$ de 16 a 20) indexam uma
tabela que aponta para o próximo salto ou para um bloco com busca binária~\cite[\S3.2]{zec12}. O Poptrie é uma trie de 64
vias em que o índice do filho sai de uma contagem de bits (\texttt{popcnt}) sobre um vetor de 64 bits, com uma tabela de
topo indexada pelos $s$ primeiros bits~\cite[\S3]{asai15}. O SAIL separa a descoberta do comprimento do prefixo, feita
com bitmaps on-chip, da busca do próximo salto, feita em arrays off-chip, e trata à parte os prefixos de até /24 e os de
/25 em diante. Na variante SAIL\_L, os prefixos são expandidos para os níveis 16, 24 e 32~\cite[\S I-C, \S IV]{yang18}.

As quatro pagam o lookup em acessos à memória e dependem de a parte mais usada caber no nível rápido, seja SRAM on-chip
ou cache. Em relação à TCAM, o custo sai da energia e vai para a hierarquia de memória e para a reconstrução de estruturas
a cada atualização.

% =====================================================================
\section{Análise por eixo}

\subsection{Velocidade de lookup}
A Tabela~\ref{tab:velocidade} reúne os números de velocidade com a condição em que foram obtidos. Três diferenças impedem a
leitura direta. A unidade varia entre lookups por segundo num laço de CPU (Mlps), pacotes por segundo (Mpps) e acessos por
lookup no pior caso. A plataforma vai de um FPGA simulado a CPUs com caches de tamanhos diferentes. E nenhum número de CPU
inclui o processamento real de pacotes: o DXR admite que o teste sintético isola o lookup e que, num roteador, outras etapas
disputariam a cache~\cite[\S4.2]{zec12}, e pela descrição dos testes o mesmo vale para o Poptrie e o SAIL.

\begin{table}[t]
\centering
\caption{Velocidade de lookup e condição de medida}
\label{tab:velocidade}
\footnotesize
\begin{tabular}{P{1.35cm} P{1.45cm} P{3.0cm} P{1.3cm}}
\toprule
Esquema & Valor & Condição & Fonte \\
\midrule
TCAM & $>$100 Mlps & dispositivo de 18 Mbit, dado de fabricante & \cite[\S II]{zane03} \\
Tree Bitmap & 25 M/s & meta de pior caso de um projeto de hardware, sem medição & \cite[\S7.2]{eath04} \\
Tree Bitmap & 56,24 Mlps & medido pelo Poptrie, 1 núcleo, aleatório & \cite[Tab.~3]{asai15} \\
DXR D16R & 490 Mlps & 8 núcleos, chaves uniformes & \cite{zec12} \\
DXR D18R & 179,92 Mlps & medido pelo Poptrie, 1 núcleo, aleatório & \cite[Tab.~3]{asai15} \\
Poptrie18 & 240,52 Mlps & 1 núcleo i7-4770K, aleatório & \cite[Tab.~3]{asai15} \\
SAIL\_L & 158,22 Mlps & medido pelo Poptrie, mesma condição & \cite[Tab.~3]{asai15} \\
SAIL\_L & 231,47--236,08 Mpps & 1 núcleo i7-3520M, aleatório & \cite[\S VI-C]{yang18} \\
SAIL\_L & 673,22--708,71 Mpps & idem, trace do tier-1 & \cite[\S VI-C]{yang18} \\
SAIL\_L & 479 Mpps & FPGA simulado, só estágios on-chip & \cite[\S VI-B]{yang18} \\
\bottomrule
\end{tabular}
\end{table}

\subsubsection{O caso Poptrie e SAIL}
O Poptrie reimplementou o SAIL e mediu o Poptrie18 de 1,37 a 2,62 vezes mais rápido com tráfego aleatório em 35
tabelas~\cite[\S4.5]{asai15} e 1,22 vez mais rápido com um trace real~\cite[\S4.7]{asai15}. O SAIL, usando o código aberto
do Poptrie, mediu o SAIL\_L de 1,13 a 1,14 vez mais rápido com tráfego aleatório e de 1,59 a 1,72 vez com tráfego
real~\cite[\S VI-C]{yang18}. Os autores do SAIL atribuem o resultado do Poptrie a uma reimplementação que
``probably introduces some unnecessary overhead''~\cite[\S VI-C]{yang18}. A Fig.~\ref{fig:razoes} põe as razões dos dois
artigos na mesma escala. Antes de tratar os números como contraditórios,
comparamos as condições de cada medida (Tabela~\ref{tab:fatores}).

\begin{figure}[t]
\centering
\begin{tikzpicture}
\begin{axis}[
  scale only axis, width=0.5\columnwidth, height=4.2cm,
  xmode=log, xmin=0.5, xmax=3,
  xtick={0.5,0.75,1,1.5,2,3},
  log ticks with fixed point,
  xticklabel style={/pgf/number format/.cd, use comma, fixed, precision=2},
  ytick={1,2,3,4,5,6,7}, y dir=reverse, ymin=0.4, ymax=7.6,
  yticklabels={Poptrie: aleatório, Poptrie: trace de borda, Poptrie: sequencial, Poptrie: repetido,
               SAIL: aleatório, SAIL: por prefixo, SAIL: trace tier-1},
  tick label style={font=\footnotesize},
  xlabel={vazão do Poptrie / vazão do SAIL\_L (escala log)},
  xlabel style={font=\footnotesize},
  xmajorgrids, grid style={gray!30},
]
\draw[dashed] (axis cs:1,0.4) -- (axis cs:1,7.6);
\draw[very thick, |-|] (axis cs:1.37,1) -- (axis cs:2.62,1);
\addplot[only marks, mark=*] coordinates {(1.22,2) (0.888,3) (0.976,4)};
\draw[very thick, |-|] (axis cs:0.877,5) -- (axis cs:0.885,5);
\draw[very thick, |-|] (axis cs:0.758,6) -- (axis cs:0.813,6);
\draw[very thick, |-|] (axis cs:0.581,7) -- (axis cs:0.629,7);
\end{axis}
\end{tikzpicture}
\caption{Razão entre a vazão do Poptrie e a do SAIL\_L em cada condição, segundo quem mediu. À direita da linha
tracejada, o Poptrie é mais rápido. Barras são intervalos entre tabelas; pontos são medidas únicas. Fontes:
\cite[\S4.5, \S4.7]{asai15} e \cite[\S VI-C]{yang18}. As razões sequencial e repetida e a inversão das razões do SAIL
são conta nossa.}
\label{fig:razoes}
\end{figure}

\begin{table}[t]
\centering
\caption{Condições das medições cruzadas entre Poptrie e SAIL}
\label{tab:fatores}
\footnotesize
\begin{tabular}{P{1.35cm} P{3.05cm} P{3.05cm}}
\toprule
Fator & Poptrie~\cite{asai15} & SAIL~\cite{yang18} \\
\midrule
CPU e cache & i7-4770K, L3 de 8 MiB, 1 núcleo (\S4) & i7-3520M, L3 de 4 MB, 1 núcleo (\S V-B, \S VI-C) \\
Rival & SAIL\_L reimplementado, sem parâmetros informados (\S4) & código aberto do Poptrie, configuração não informada (\S VI-C) \\
Demais concorrentes & todos reimplementados (\S4) & 4 de 5 implementados pelos autores (\S VI-A) \\
Tráfego aleatório & xorshift, custo do gerador incluído (\S4.2) & geração não detalhada (\S VI-C) \\
Tráfego real & roteador de borda de AS; 21,8\% dos pacotes casam $>$/24 (\S4.2, \S4.7) & roteador tier-1 na China; fração $>$/24 não informada (\S VI-A) \\
Memória do SAIL & 44,24 MiB no total (Tab.~3) & no máximo 2,13 MB on-chip (\S IV-D) \\
\bottomrule
\end{tabular}
\end{table}

Quatro fatores explicam a divergência sem que nenhum dos lados tenha errado as contas. O primeiro é o tráfego. O SAIL
declara que o seu escopo é o backbone, onde quase todo pacote casa um prefixo de até /24, e admite que em roteadores de
borda mais tráfego cai em prefixos longos e o lookup fica mais lento~\cite[\S I-C, \S VI-A]{yang18}. O trace real do Poptrie
vem de um roteador de borda de AS, e nele 21,8\% dos pacotes casam prefixos maiores que /24, contra 1,66\% se os destinos
fossem uniformes no espaço IPv4~\cite[\S4.2, \S4.7]{asai15}. Os dois traces reais estão, portanto, em lados opostos da
premissa que o próprio SAIL declara. Há um ponto fraco do lado do SAIL: a premissa não é quantificada no trace do tier-1.
O único número de tráfego acima de /24 que o artigo apresenta é para tráfego aleatório, de 0,021\% a
0,025\%~\cite[\S VI-C, Fig.~11]{yang18}.

Depois vem a combinação de localidade e cache. Nos dados do próprio Poptrie, o SAIL vence com tráfego sequencial
(1264 contra 1122~Mlps) e repetido (492 contra 480~Mlps)~\cite[\S4.5]{asai15}. A distribuição de ciclos por lookup ajuda a
entender: o SAIL tem mediana de 22 ciclos e percentil 99 de 299, e o Poptrie18 tem 46 e 169~\cite[Tab.~4]{asai15}. Pela
nossa leitura, o SAIL é muito rápido quando o dado está na cache e lento quando precisa ir à DRAM. Com 44,24~MiB de
estrutura e 8~MiB de L3, o tráfego aleatório leva o SAIL ao segundo caso, e o tráfego com localidade, como o do trace do
tier-1, ao primeiro.

Cada lado também mediu o rival com uma implementação que não controlava por completo. O SAIL
mede a si mesmo a 231,47--236,08~Mpps com tráfego aleatório, numa CPU com metade da L3~\cite[\S VI-C]{yang18}, e o
Poptrie mede o SAIL a 158,22~Mlps~\cite[Tab.~3]{asai15}. Essa diferença é compatível com a queixa do SAIL sobre a
reimplementação. O SAIL, porém, tem o mesmo viés: implementou ele próprio quatro dos cinco concorrentes de
CPU~\cite[\S VI-A]{yang18} e não informa a configuração do Poptrie que usou.

Por último, a forma de medir. O Poptrie inclui na taxa o custo do gerador de endereços, 1,22~ns por
lookup~\cite[\S4.2]{asai15}. A 240,52~Mlps, cada lookup leva 4,16~ns, dos quais cerca de 29\% são do gerador (conta
nossa), o que reduz as razões entre os algoritmos. O SAIL não descreve como gera o tráfego aleatório.

Os dois resultados, portanto, delimitam cenários. Com tráfego de backbone concentrado em prefixos de até /24 e boa
localidade, o SAIL\_L lidera nas duas medições que se aproximam dessa condição: o tráfego sequencial no Poptrie e o trace do
tier-1 no SAIL. Com tráfego aleatório, ou com muitos pacotes em prefixos longos, como numa borda com rotas de IGP mais
específicas, o Poptrie lidera. O SAIL também testa FIBs com 1\% a 10\% de prefixos longos e continua à frente, por 1,47 a
2,83 vezes~\cite[\S VI-C, Fig.~14]{yang18}. Esse experimento, porém, varia a fração de prefixos longos na tabela, e a
disputa com o Poptrie depende da fração do tráfego que chega a eles. \revisar{conferir na Fig.~14 do yang18 qual tráfego foi usado}

\subsubsection{DXR, Tree Bitmap e TCAM}
O DXR chama o tráfego uniforme aleatório de pior caso~\cite[\S4.1]{zec12}. Para uma estrutura de índice direto, esse
padrão é favorável: de 75,7\% (D16R) a 94,5\% (D20R) das chaves se resolvem já no índice, sem busca
binária~\cite[Tab.~3]{zec12}. Nas medições do Poptrie, a vantagem dele sobre o D18R é pequena, de 1,04 a 1,34
vez~\cite[\S4.5]{asai15}. Por outro lado, o percentil 95 do DXR e do SAIL passa de 234 ciclos nas profundidades 24 e
25~\cite[\S4.6]{asai15}, o que no DXR bate com o custo da busca binária em prefixos longos. O Tree Bitmap não traz nenhuma
medição de taxa: a sua Tabela~5 é analítica, e a implementação em hardware é um projeto com meta de 25 milhões de lookups
por segundo no pior caso~\cite[\S7.2]{eath04}. Medido em software pelo Poptrie, fica de 3,52 a 6,78 vezes
atrás~\cite[\S4.5]{asai15}, comparação que não faz justiça a um esquema pensado também para hardware. Para a TCAM, o
corpus só tem números de fabricante, ``over 100 Million lookups/second''~\cite[\S II]{zane03}. Nenhum artigo algorítmico
mede uma TCAM real.

\subsection{Memória}
Os artigos usam ``memória'' para grandezas diferentes (Tabela~\ref{tab:memoria}). As medidas do DXR e do Poptrie contam
só a estrutura de lookup, mas os dois dependem de uma árvore radix separada para aplicar as
atualizações~\cite[\S3.4]{zec12}, \cite[\S3.5]{asai15}, e essa árvore fica fora da conta. A ordem de grandeza do que fica
de fora aparece na Tabela~3 do próprio Poptrie, em que a radix, usada como estrutura de lookup, ocupa
30,48~MiB~\cite[Tab.~3]{asai15}. No SAIL, o tamanho constante vale só para a memória on-chip, que é de no máximo 2,13~MB
qualquer que seja a FIB~\cite[\S IV-D]{yang18}. A memória off-chip cresce com a tabela; o artigo não informa o seu total em
IPv4 e pede pelo menos 4~GB de DRAM~\cite[\S IV-D]{yang18}. A TCAM é o único caso em que a tabela inteira fica no dispositivo
de lookup. A ordenação, porém, depende de um software de gerência com uma trie auxiliar fora da TCAM~\cite[p.~38]{shah01},
cujo tamanho o artigo não informa.

\begin{table}[t]
\centering
\caption{O que cada artigo chama de memória}
\label{tab:memoria}
\footnotesize
\begin{tabular}{P{1.25cm} P{2.65cm} P{2.4cm} P{1.0cm}}
\toprule
Esquema & O que é contado e valor & Fica de fora & Fonte \\
\midrule
TCAM & entradas de largura fixa; 18 Mbit $\approx$ 512 mil prefixos IPv4 & trie auxiliar de gerência & \cite[\S II]{zane03} \\
Tree Bitmap & estrutura de busca: 11--23 B/prefixo; físico com réplicas: 128 B/prefixo & folhas, na 1.\textsuperscript{a} medida & \cite[Tab.~6, \S7.2]{eath04} \\
DXR D16R & estrutura de lookup: 782 KB para 417 mil prefixos & radix usada para reconstruir & \cite[\S3.4]{zec12} \\
Poptrie18 & FIB: 2,40 MiB & RIB em radix (a radix ocupa 30,48 MiB) & \cite[Tab.~3, \S3]{asai15} \\
SAIL\_L & limite on-chip: 2,13 MB para qualquer FIB & arrays off-chip; 44,24 MiB no total, medido pelo Poptrie & \cite[\S IV-D]{yang18}, \cite[Tab.~3]{asai15} \\
\bottomrule
\end{tabular}
\end{table}

Dois números do Tree Bitmap pedem cautela. O artigo informa 128~B de DRAM física por prefixo e 1,6~GB para 10 milhões de
prefixos~\cite[\S7.2]{eath04}, mas 128~B vezes 10 milhões dá 1,28~GB, e a diferença não é explicada.
\revisar{conferir a conta na p.~120 do eath04} O argumento de custo do mesmo artigo é mais sólido: por bit, a SRAM custa até 36
vezes o preço da DRAM de rede com a mesma banda~\cite[\S2.3.2]{eath04}, o que explica a escolha de DRAM no projeto.

Para a parte que precisa ficar no nível rápido, DXR e Poptrie cabem com folga numa L3 de 8~MiB, o que atende roteadores em
software. Em hardware com SRAM on-chip limitada e FIB crescente, o limite fixo do SAIL é a única garantia do corpus que não
depende do tamanho da tabela. Para chaves largas e muitos prefixos, o custo por bit favorece a DRAM e pesa contra a TCAM.

\subsection{Custo de atualização}
O requisito mais citado é de Labovitz, que chega até nós por fonte secundária: em roteadores de backbone, ocorrem em média
100 mudanças de rota por segundo, com rajadas acima de ``several hundred'' por segundo. Por isso Ruiz-Sánchez et al.
exigem no máximo 10~ms por atualização~\cite[p.~11]{ruiz01}. Shah e Gupta citam outro trabalho do mesmo grupo para 100 a
1000 prefixos por segundo~\cite[p.~38]{shah01}, e o CoolCAMs observa em traces reais rajadas de até ``a few thousand''
adições de rota num único segundo~\cite[\S V]{zane03}. Os 10~ms saem da média. Numa rajada de alguns milhares por segundo,
o orçamento cai para menos de 1~ms por atualização (conta nossa). A Tabela~\ref{tab:update} compara o que cada artigo
mede.

\begin{table*}[t]
\centering
\caption{Custo de atualização: unidade, valores e o que fica fora da conta}
\label{tab:update}
\footnotesize
\begin{tabular}{P{2.4cm} P{2.3cm} P{2.6cm} P{3.4cm} P{3.0cm} P{2.6cm}}
\toprule
Esquema & Unidade & Média & Pior caso & Fora da conta & Carga de teste \\
\midrule
TCAM, CAO\_OPT~\cite{shah01} & movimentos de memória & 1,02--1,05 (Tabs.~2--3) & 3 observado; $D/2$ teórico, otimalidade só conjecturada (pp.~40, 43) & tempo, CPU da trie auxiliar, lookups durante o update & traços de 3 h; tabelas de 35--43 mil (Tab.~1) \\
TCAM particionada~\cite{zane03} & escritas na TCAM e recomputações & 1 a 1,07 escrita; 0 a 74 recomputações com 8 buckets (Tabs.~III--IV) & reescrever a TCAM de dados inteira (\S V) & movimentos para manter a ordem, delegados a~\cite{shah01} & 3,4 e 3,6 milhões de updates (Tab.~III) \\
Tree Bitmap~\cite{eath04} & cópias de nó & não medido & $256+C$, afirmado sem derivação (\S4) & expansão no array inicial (\S5.1) & nenhuma \\
DXR~\cite{zec12} & \textmu s por prefixo & 10,2--11,9\us{} (\S4.4) & não informado; prefixos curtos custam mais (Fig.~7) & atualização da radix (\S4.4) & cada prefixo da LINX removido e reinserido (\S4.4) \\
Poptrie18~\cite{asai15} & \textmu s por update & 2,51\us{} (\S4.9) & não medido; até 1~MiB copiado ao trocar o array de topo (\S3.4--3.5, conta nossa) & --- & 23.446 updates de 1 h, um único peer (\S4.9) \\
SAIL\_L~\cite{yang18} & acessos à memória & 1,807--2,253 (Fig.~12) & 7,88 observado; ilimitado em teoria (Tab.~II) & trie do plano de controle (\S IV-B) & $319\times500$ updates de 3 coletores (\S VI-C) \\
\bottomrule
\end{tabular}
\end{table*}

Na média, todas as soluções algorítmicas medidas ficam de duas a três ordens de grandeza abaixo dos 10~ms, e o CAO\_OPT faz
pouco mais de um movimento por atualização. Pelo critério de 2001, o eixo estaria resolvido, mas nenhum artigo mede o pior
caso nem o comportamento em rajadas. O SAIL\_L tem pior caso ilimitado~\cite[Tab.~II, \S IV-E]{yang18}, e enquanto a
estrutura on-chip muda nenhum lookup pode ser feito, o que pode causar perda de pacotes~\cite[\S III-C]{yang18}. O Poptrie
atualiza sem bloquear o lookup, com um único escritor e troca atômica de ponteiro~\cite[\S3.5]{asai15}, mas só informa a
média. Quando o nó afetado fica acima da profundidade $s$, o array de topo inteiro, de até $4\times2^s$ bytes, é
substituído~\cite[\S3.4--3.5]{asai15}; com $s=18$, isso dá 1~MiB (conta nossa). O DXR deixa fora da medida a atualização da
radix~\cite[\S4.4]{zec12}. O limite do Tree Bitmap é afirmado sem derivação, com uma constante $C$ que o texto não define, e
não cobre a expansão no array inicial. O próprio artigo admite que arrays iniciais de 13 bits ou mais podem ter
``poor update times'', e 13 bits é o valor da implementação de referência em
software~\cite[\S4, \S5.1, \S6]{eath04}.

Do lado da TCAM, o pior caso do CAO\_OPT, $D/2$, depende do comprimento das cadeias de prefixos aninhados, que nas tabelas
medidas não passava de 5~\cite[p.~40]{shah01}. A vantagem depende, portanto, das cadeias curtas dessas tabelas, já que o
limite cresce com $D$. As tabelas tinham de 35 a 43 mil prefixos, uma ordem de grandeza a menos que as do DXR e do Poptrie, e os traços
de 3~h correspondem a cerca de 3 a 4 atualizações por segundo (conta nossa a partir da Tab.~1), longe das rajadas que
motivam o artigo. \revisar{na Tab.~1 do shah01, ``Prefixes'' é igual a Inserts + Deletes; confirmar na p.~43 antes de usar o
tamanho das tabelas} O CoolCAMs conta uma escrita na TCAM por atualização e delega a manutenção da ordem ao trabalho de
Shah e Gupta~\cite[\S II]{zane03}, de modo que os dois custos nunca foram medidos juntos.

Se o critério for a média, DXR e Poptrie vencem em software, e o CAO\_OPT deixa a TCAM competitiva em número de
movimentos. Se o critério for o pior caso com lookups ocorrendo ao mesmo tempo, o corpus não aponta um vencedor. O único
limite pequeno garantido pela própria estrutura é o do SAIL\_U, de um acesso on-chip por atualização, ao custo de quatro
acessos on-chip por lookup em vez dos dois do SAIL\_L~\cite[Tab.~II]{yang18}.

\subsection{Energia e custo}
Os três artigos algorítmicos recentes justificam o abandono da TCAM pelo consumo: ``power-hungry
TCAMs''~\cite[\S1]{zec12}, problemas de energia e calor~\cite[\S1]{asai15}, ``a huge amount of power''~\cite[\S II]{yang18}.
Nenhum deles mede o próprio consumo nem apresenta um número de TCAM.

Os dois artigos que discutem a TCAM com números chegam a conclusões opostas. O CoolCAMs a chama de
``cost-effective''~\cite{zane03} e mostra que, se a tabela for particionada e só um bucket for pesquisado, a potência
cai de forma garantida em 7,98 vezes com 8 buckets~\cite[\S IV-D]{zane03}. O Tree Bitmap considera o custo
``extremely prohibitive'' e estima que 10 milhões de entradas de 52 bits exigiriam 20 componentes, mais de US\$~4000 e
mais de 200~W~\cite[\S2.3.3]{eath04}. Sobre o CoolCAMs, diz que as ideias dele ``do not appear to be reflected in
practice''~\cite[\S3]{eath04}.

Os cenários, porém, são diferentes. O CoolCAMs trata de tabelas IPv4 de até 1 milhão de entradas~\cite[\S II]{zane03}. O
Tree Bitmap trata de um roteador de múltiplos serviços que pode precisar de chaves de VPN de 52 bits e de até 10 milhões de
prefixos~\cite[\S2.1, Tab.~3]{eath04}, cenário que o próprio artigo diz não ser o comum~\cite[\S7.2]{eath04}. Além disso,
os 200~W só valem ``if techniques are not employed to limit the components and banks of CAMs that are
searched''~\cite[\S2.3.3]{eath04}, e limitar os bancos pesquisados é a técnica do CoolCAMs. As duas posições são
compatíveis. Com particionamento e uma tabela IPv4 de 1 milhão de entradas, a TCAM fica em poucos watts segundo o modelo; com
chaves largas, dezenas de chips e busca em todos os bancos, fica cara.

Os números dos dois lados são frágeis. O CoolCAMs não mede potência: multiplica o número de entradas pesquisadas pela
potência nominal de 15~W, com um modelo linear sem citação de fonte~\cite[\S II, \S IV-D]{zane03}. O fator 7,98 vale para
uma das variantes (\textit{postorder-split}) numa tabela hipotética de exatamente 1 milhão de entradas, e as tabelas reais
testadas tinham de 110 a 122 mil prefixos~\cite[Tab.~II]{zane03}. O exemplo de pior caso usa uma TCAM de dados de 15~W e
1 milhão de entradas~\cite[\S IV-D]{zane03}, enquanto a Seção~II do artigo associa 15~W a 512 mil prefixos. Pelo modelo do
próprio artigo, pesquisar 125.288 entradas daria cerca de 3,6~W, e não menos de 2~W (conta nossa).
\revisar{conferir a conta} No Tree Bitmap, com chips de 256 mil entradas de 72 bits~\cite[\S2.3.3]{eath04}, 10 milhões de
entradas pediriam cerca de 39 componentes. O número 20 só fecha com entradas de 36 bits, em que 52 bits não cabem (conta
nossa). \revisar{conferir a conta} O erro subestima a TCAM e não muda a conclusão do artigo, mas o número publicado não é
reprodutível.

Com o corpus, não dá para dizer qual abordagem gasta menos energia por lookup, porque só a TCAM tem modelo de consumo e
ninguém mede o das soluções algorítmicas. O corpus sustenta afirmações mais modestas: o particionamento reduz o consumo
modelado da TCAM em cerca de 8 vezes para tabelas IPv4 de 1 milhão de entradas; chaves largas e tabelas de dezenas de
milhões encarecem a TCAM; e, para funções de rede virtualizadas em servidores comuns, não há TCAM disponível~\cite[\S1]{asai15}, o que resolve
a questão sem depender de números.

\subsection{Escala da tabela e IPv6}
O survey inclui entre os critérios a escalabilidade com o número e o comprimento dos prefixos~\cite[p.~20]{ruiz01}. Os
artigos do corpus testam escalas muito diferentes (Tabela~\ref{tab:corpus}). Os dois de TCAM usam tabelas reais de 35 a
122 mil prefixos, e o caso de 1 milhão de entradas do CoolCAMs aparece só na análise de pior caso e numa tabela
sintética~\cite[\S III-D, \S IV-D]{zane03}. O Tree Bitmap avalia tabelas de 19 a 41 mil prefixos~\cite[Tab.~6]{eath04},
embora fixe para a tabela global um requisito de ``a million or more prefixes''~\cite[\S2.2]{eath04}. DXR, Poptrie e SAIL
usam tabelas reais de 114 a 532 mil prefixos, e só o Poptrie vai além, com tabelas sintéticas de até 886
mil~\cite[Tab.~1]{asai15}.

Nessa faixa aparecem limites de formato que os testes com tabelas menores não mostram. O SAIL guarda o identificador de
chunk em 15 bits, o que bastava nas FIBs reais medidas pelos autores~\cite[\S IV-B]{yang18}. Nas duas tabelas sintéticas
maiores do Poptrie, de 877 e 886 mil prefixos, o número de chunks passa de $2^{15}$ e o SAIL não
compila~\cite[\S4.8, Tab.~5]{asai15}. A versão de 2018 do SAIL chama a falha de ``tiny problem'' e propõe 16 bits mais um
bitmap de $2^{16}$ bits~\cite[\S VI-C]{yang18}, sem medir a correção. O limite depende do número de chunks e não do total
de rotas, e nenhum dos dois artigos dá um limiar em número de rotas. Essas tabelas sintéticas foram construídas dividindo
cada prefixo de até /24, /20 e /16 em dois, quatro e oito pedaços com próximos saltos distintos~\cite[\S4.1]{asai15}, o
que multiplica os chunks. Pela nossa leitura, não se sabe se tabelas reais do mesmo tamanho teriam o mesmo efeito. O DXR
tem limites do mesmo tipo, até $2^{19}$ faixas e 4096 entradas por chunk~\cite[\S3.2]{zec12}, e o Poptrie precisou
estender o índice do DXR para $2^{20}$ faixas para rodá-lo nessas tabelas~\cite[\S4.8]{asai15}. O próprio Poptrie limita
a FIB a $2^{16}$ entradas, porque a folha tem 16 bits, e diz que o campo pode crescer ao custo de mais
memória~\cite[\S5]{asai15}.

No IPv6, a evidência é fraca em todas as famílias. O DXR declara o índice direto inviável, porque o espaço IPv6 é
esparso: 0,026\% de ocupação, contra 60,36\% no IPv4~\cite[\S3]{zec12}. Mesmo assim, o Poptrie adaptou o D18R ao IPv6 e o
mediu a 169,91~Mlps~\cite[\S4.10]{asai15}, numa tabela pequena e com endereços sorteados dentro de 2000::/8, o que não
decide a questão. Segundo o Poptrie, o SAIL não aceita prefixos mais longos que /64~\cite[\S4.10]{asai15}; a versão de
2018 propõe um híbrido de bitmaps e filtros de Bloom, com até 12 acessos off-chip no pior caso, avaliado em FIBs IPv6
sintéticas derivadas das IPv4~\cite[\S VII-C]{yang18}. O Poptrie é o único com teste numa tabela IPv6 real, de 20.440
prefixos, em que o Poptrie18 faz 211,32~Mlps, 1,24 vez o D18R~\cite[\S4.10, Tab.~6]{asai15}; os próprios autores chamam
essa avaliação de ``less interesting'' por causa do tamanho da tabela~\cite[\S4.10]{asai15}. Do lado da TCAM, o CoolCAMs
diz que técnicas semelhantes servem ao IPv6, sem avaliá-las~\cite[\S I]{zane03}, e Shah e Gupta só definem
$L=128$~\cite[p.~37]{shah01}. Pela fórmula, o pior caso do PLO\_OPT passaria de 16 para 64 movimentos (conta nossa).

Nenhuma abordagem tem evidência para tabelas reais maiores que 532 mil prefixos ou para IPv6 em escala de produção. Nas
tabelas sintéticas maiores, os limites de formato apareceram antes dos de desempenho, inclusive no SAIL, que se apresenta
como independente do tamanho da FIB.

% =====================================================================
\section{Discussão: onde cada abordagem vence}

A Tabela~\ref{tab:cenarios} junta as análises da Seção~IV. Para cada cenário, indicamos a abordagem mais bem sustentada
pelo corpus e o tipo de evidência: \textbf{M} quando os próprios autores mediram, \textbf{C} quando um concorrente mediu e
\textbf{A} quando há só análise, modelo ou simulação sem tráfego. Nenhum resultado do corpus foi reproduzido por um grupo
independente; as medições cruzadas entre Poptrie, SAIL e DXR foram feitas por concorrentes diretos.

\begin{table*}[t]
\centering
\caption{Abordagem mais bem sustentada por cenário (M: medido pelos autores; C: medido por concorrente; A: análise ou modelo)}
\label{tab:cenarios}
\footnotesize
\begin{tabular}{P{4.3cm} P{2.6cm} P{3.4cm} P{5.6cm}}
\toprule
Cenário & Abordagem & Evidência & Ressalva principal \\
\midrule
Backbone em hardware com SRAM on-chip (ASIC ou FPGA), FIB crescendo, tráfego em prefixos de até /24 & SAIL\_L & M (CPU, FPGA simulado); A (limite on-chip) & pior caso de atualização ilimitado; lookup bloqueado durante a escrita on-chip (\S III-C, Tab.~II de~\cite{yang18}) \\
Roteador em software ou NFV em x86, tráfego de borda com prefixos longos, IPv4 e IPv6 & Poptrie & M (35 tabelas, trace de borda); C (atrás do SAIL com tráfego de backbone) & laço de lookup sem E/S de pacotes; IPv6 testado com 20.440 prefixos (\S4.10 de~\cite{asai15}) \\
Software, só IPv4, menor pegada de memória & DXR (D16R) & M; C (Poptrie e SAIL) & só IPv4 (\S3 de~\cite{zec12}); busca binária lenta em prefixos longos \\
Hardware com DRAM, chaves largas (VPN, várias tabelas), pior caso determinístico & Tree Bitmap & A (sem medição no artigo) & limite de atualização sem prova; tabelas de 19--41 mil \\
Tabelas pequenas ou médias com exigência de determinismo e gerência simples & TCAM com CAO\_OPT e particionamento & A (simulação de updates e modelo de potência) & custo e potência crescem com o número de entradas; tabelas testadas de 35 a 122 mil \\
Tabelas grandes com poucos pacotes acima de /24 & híbrido: SAIL para até /24, TCAM para /25 ou mais & proposto, não medido (\S I-D de~\cite{yang18}) & sem avaliação no corpus \\
\bottomrule
\end{tabular}
\end{table*}

Entre as soluções algorítmicas, a escolha depende mais da composição do tráfego e do tamanho da cache que do tamanho da
tabela. O SAIL e o Poptrie trocam de posição conforme a fração de pacotes que alcança prefixos longos, e o DXR perde
terreno nas profundidades em que a busca binária domina. Para a TCAM, pesam o tamanho da tabela e a largura da chave: o
CoolCAMs e o Tree Bitmap concordam nos fatos e divergem porque tratam de pontos diferentes desse
espaço.

Os cenários com IPv6 e com tabelas reais acima de 532 mil prefixos ficaram fora da tabela, porque nenhuma abordagem tem
evidência suficiente para eles (Seção~IV-E).

Os híbridos aparecem em três artigos, cada um de um jeito. O SAIL admite TCAM para os prefixos de /25 em
diante~\cite[\S I-D]{yang18}, e o CoolCAMs usa uma trie numa TCAM de índice pequena para escolher o
bucket~\cite[\S IV]{zane03}. Os ``nós CAM'' do Tree Bitmap, por sua vez, são uma codificação de nó em DRAM, não um chip de
TCAM~\cite[\S5.5]{eath04}. Nenhum desses híbridos tem medição de desempenho no corpus.

As divergências da Seção~IV indicam o que uma comparação entre essas estruturas precisaria informar para ser
aproveitada por outros grupos. Cada item abaixo responde a um problema que encontramos:
\begin{enumerate}
\item a origem do trace (borda ou backbone) e a fração dos pacotes que casa prefixos acima de /24, da qual depende a
  disputa entre Poptrie e SAIL;
\item a origem do código de cada concorrente, com versão e parâmetros, ou a reimplementação publicada;
\item a distribuição de ciclos por lookup, além da taxa média, com o custo de gerar os endereços descontado;
\item a memória por nível (on-chip ou cache, e off-chip), incluindo a estrutura de controle usada nas atualizações;
\item o pior caso de atualização e o comportamento em rajadas de milhares de atualizações por segundo, com lookups
  acontecendo ao mesmo tempo;
\item a energia por lookup, medida da mesma forma para a TCAM e para as estruturas algorítmicas.
\end{enumerate}
O survey de 2001 já fazia parte disso: publicava percentis do tempo de lookup e tratava o tráfego uniforme por prefixo
como limitação~\cite[Tab.~3, p.~22]{ruiz01}. Entre os três artigos algorítmicos recentes, só o Poptrie publica a
distribuição de ciclos por lookup~\cite[Tab.~4]{asai15}. O DXR dá a distribuição de passos da busca
binária~\cite[Tab.~3]{zec12}, e no SAIL não encontramos distribuição de latência.

% =====================================================================
\section{Limitações e ameaças à validade}

O corpus termina em 2018 e tem um representante por família. Ficaram de fora trabalhos mais recentes e plataformas como
switches programáveis, e os dados de capacidade e potência de TCAM são de dispositivos de 2001 a 2004. As conclusões sobre
famílias valem, portanto, para os representantes escolhidos.

Quase toda a evidência foi medida pelos próprios autores, e as medições cruzadas foram feitas por concorrentes. Encontramos
também inconsistências internas que reduzem a reprodutibilidade. No SAIL, há ``12 FIBs'' listadas com 13
nomes~\cite[\S VI-A]{yang18}, uma tabela anunciada com 23 FIBs que tem 24 linhas~\cite[Tab.~IV]{yang18}, três tamanhos
diferentes para a mesma tabela rrc00~\cite[\S VI-B, \S VII-C, Tab.~IV]{yang18} e citações trocadas, como ``Poptrie
[11]''~\cite[\S VI-C]{yang18}. No Tree Bitmap, os doze valores da Tabela~5 ficam 9 bits acima da fórmula dada no próprio
texto~\cite[\S7.2]{eath04}. Evitamos usar números que dependem desses pontos.

As contas marcadas como ``conta nossa'' são derivações simples a partir de números impressos, e não usamos leituras
aproximadas de gráficos. Por fim, o primeiro rascunho foi escrito por IA. Os riscos conhecidos desse uso são erros
plausíveis e números trocados; a Seção~II-C descreve como eles foram controlados.
\revisar{completar com o resultado da revisão do grupo}

% =====================================================================
\section{Conclusão}

Os resultados publicados só permitem escolher entre TCAM e algoritmo depois de fixado o cenário. Em roteadores de software, Poptrie e DXR cabem na cache e atualizam em microssegundos. O Poptrie leva vantagem com
prefixos longos e IPv6, e o DXR com a menor pegada em IPv4. Em hardware de backbone com SRAM on-chip, o SAIL é o único a
garantir memória on-chip constante, ao preço de um pior caso de atualização ilimitado. A TCAM continua defensável para
tabelas de tamanho moderado, em que determinismo e simplicidade pesam mais que energia, e o particionamento do CoolCAMs
reduz o consumo modelado em cerca de 8 vezes.

Nenhum dos sete artigos traz as medições que mais ajudariam a decidir: o pior caso de atualização durante rajadas, com
lookups acontecendo ao mesmo tempo, em todas as estruturas; a energia por lookup das soluções algorítmicas, na mesma base
usada para a TCAM; e testes com tabelas reais maiores, já que nas sintéticas os limites de formato apareceram. Também faltam
traces públicos de borda e de backbone que informem a fração de pacotes acima de /24, a premissa de que depende a disputa
entre Poptrie e SAIL.

\bibliographystyle{IEEEtran}
\bibliography{refs}

\end{document}
